\printedpage{39}
% head: A. Grothendieck : Dérivateurs, Ch. VII, Édition des Universités de Montpellier II et Paris VII
On the other hand, setting $l_i = \mathrm{long}(\pi_{i-1})$ (for $i \in [1,r]$),
\[
f_i : \Delta_\varepsilon^{l_{i-1}} \longrightarrow \Delta_\varepsilon^{l_i\ (=\mathrm{long}(\pi_{i-1}))}
\]
is defined as
\[
f_i = \mathfrak{f}_{\varepsilon,\pi_{i-1}(1)} * \mathfrak{f}_{\varepsilon,\pi_{i-1}(2)} * \cdots * \mathfrak{f}_{\varepsilon,\pi_{i-1}(l_{i-1})},
\]
so that the element $(*)$ of $\Nerf_r(\Ty)$ is uniquely determined by the data of the $\varepsilon \in \{\pm 1\}$ and of
\[
\pi_\bullet = (\pi_i)_{i\in[0,r-1]} \in (\mathrm{Hom}_{\mathrm{Ens}}(\mathbf{N}^*,\mathbf{N}))^r
\]
satisfying the conditions $(**)$. Let $F_{\varepsilon,\pi_\bullet}$ be this element of $\Nerf_r(\Ty)$, and let $N_r$ be the set of $\pi_\bullet$ we have just made explicit. Then
\[
(F_{\varepsilon,\pi_\bullet})_{(\varepsilon,\pi_\bullet)\in\{\pm1\}\times N_r}
\]
is a free system of generators of the monoid $\Nerf_r(\Ty)$.

If $l_i = 0$ (i.e. \fbox{$\pi_{i-1} = 0$}), we must have $\mathrm{supp}(\pi_{i+1}) = \emptyset$, i.e. $\pi_{i+1} = 0$, whence $l_{i+1} = 0$, and likewise for the following $l_j$ and $\pi_j$. Geometrically the hypothesis is expressed by the condition that the target
\msp{page 58}
of $f_i$ is the element 1 of Ty (corresponding to the one-point interval), and since $\mathrm{Hom}(1,\tau) = \emptyset$ for $\tau \neq 1$, this means that the arrows following $f_i$ must all be equal to the arrow $\mathrm{id}_1 : 1 \longrightarrow 1$.

\textbf{Remark.} In the combinatorial description of $\Ty$ and of $\Nerf_\bullet(\Ty)$, the set of ``signatures'' $\{\pm1\}$ plays a rather phony role. It can be replaced by any set $\Sigma$ (called the ``set of signatures''), and we then construct a monoidal category (or a categorical monoid) having as its set of objects the free monoid generated by $\Sigma$, interpreted as intervals of integers $[0,n]$, equipped with ``signatures'' in $\Sigma$ for each transition interval, or what amounts to the same thing, $\Sigma$-weighted directed ``graphs'' of combinatorial type an interval $\boldsymbol{\Delta}_i$, which amounts to giving an (ordered) sequence of $n$ elements $\varepsilon_1,\ \varepsilon_2,\ \ldots\ ,\ \varepsilon_n \in \Sigma$, i.e. a ``word'' $\varepsilon_1\varepsilon_2\ldots\varepsilon_n$, which we prefer to write $\Delta_{\varepsilon_1}\Delta_{\varepsilon_2}\cdots\Delta_{\varepsilon_n}$. We can define in the obvious way the \textit{increasing} homomorphisms of such weighted directed graphs (understood: with the source element and the target element in the structure, to be respected by the morphisms), without needing a categorical interpretation. The functoriality of the operation of composition of weighted intervals, for this notion of morphism, is obvious. The
\msp{page 59}
choice $\Sigma = \{\pm1\}$ is to my knowledge the only useful one, because of its link with the notion of paths and the calculation of $\pi_0$ of an object of $X$ \texttt{[rather of an object} $X$ \texttt{of} Cat\texttt{]}, etc.
% foot: Version 22 décembre 2001

\printedpage{40}
% head: A. Grothendieck : Dérivateurs, Ch. VII, Édition des Universités de Montpellier II et Paris VII
\section*{7\quad Explicit description of the categories of fractions $X\Sigma^{-1}$ in terms of the $\Ch(X;a,b)$}

If $X$ is a category, we denote by $\Ch(X)$ the category of paths in $X$ (cf. §4), and for $a,b \in \mathrm{Ob}\,X$, by
\[
\Ch(X;a,b) \qquad \text{or simply } \Ch(a,b)
\]
that of paths from $a$ to $b$, i.e. with source $a$, target $b$. The geometric intuition (which we will make explicit) is that such a path $c$ defines a ``direct'' path from $a$ to $b$ (i.e. a path of type $\tau'$ ``direct'', i.e. all the transition arrows are direct, i.e. $\tau'$ of the form $\Delta_+^n$), of the same length as the initial path $c$, obtained by inverting all the retrograde arrows in the path $c$. This operation is actually possible in $X$ itself if and only if the said arrows are invertible. We then obtain a composable sequence of arrows, whose total composite (which we can call ``the useful effect'' of the path)
\msp{page 60}
is an arrow from $a$ to $b$. In the general case, paths from $a$ to $b$ can be seen as a sort of ``generalized arrows'' going from $a$ to $b$, whose composition would be the obvious composition, and the identities the ``empty paths'' relative to the $a \in \mathrm{Ob}\,X$. We do moreover find in this way a category, having the same objects as $X$, but where the
\[
\mathrm{Hom}_{\mathrm{Ch}}(a,b) = \mathrm{Ch}(a,b) = \mathrm{Ob}\,\Ch(a,b)
\]
are prohibitively large. This is expressed in particular by the fact that the only invertible arrows in $X_{\mathrm{Ch}}$ are the identities! We have no natural functor between $X_{\mathrm{Ch}}$ and $X$ that is the identity on objects. For to a path $a \xrightarrow{c} b$ we cannot associate any $X$-morphisms $a \longrightarrow b$ (unless $c$ \texttt{[is]} of type $\Delta_+$), and to an $X$-morphism $a \xrightarrow{u} b$ we can associate a path $c(u)$ of type $\Delta_+$, but if $a \xrightarrow{u} b \xrightarrow{v} c$, $c(vu)$ is still of type $\Delta_+$ whereas $c(v)c(u)$ is of type $\Delta_+^2$, they are not equal. On the other hand, there is a natural arrow between these two elements of $\Ch(a,c)$, an arrow over the unique homomorphism
\[
\Delta_+ \xrightarrow{\mathfrak{f}_{+,2}} \Delta_+^2.
\]
\msp{page 61}
We then have
\[
\mathfrak{f}_{+,2}^*(\underbrace{a \xrightarrow{u} b \xrightarrow{v} c}_{c(v)c(u)}) = (\underbrace{a \xrightarrow{vu} c}_{c(vu)}),
\]
whence a canonical arrow
\[
c(vu) \longrightarrow c(v)c(u) \qquad \text{in } \Ch(a,c).
\]
This shows that $c(vu)$ and $c(v)c(u)$ are at least in the same connected component of $\Ch(a,c)$, or (as we also say) that they are ``homotopic''. This suggests here modifying the definition of $\mathrm{Hom}_{\mathrm{Ch}}(a,b)$ to
\[
\pi(a,b) = \pi_0(\Ch(a,b)).
\]
% foot: Version 22 décembre 2001

\printedpage{41}
% head: A. Grothendieck : Dérivateurs, Ch. VII, Édition des Universités de Montpellier II et Paris VII
The path composition functor
\[
\begin{array}{ccccc}
\Ch(a,b) & \times & \Ch(b,c) & \longrightarrow & \Ch(a,c)\\
(c & , & c') & \longmapsto & c' \circ c
\end{array}
\]
\texttt{[there are too many} $c$\texttt{]} (\footnote[17]{This functor is deduced from the functors
\[
\Ch_\tau(a,b) \times \Ch_{\tau'}(b,c) \longrightarrow \Ch_{\tau\tau'}(a,c)
\]
which are obvious, cf. page 77.}) induces a map on the $\pi_0$, which gives a composition law
\[
\pi(a,b) \times \pi(b,c) \longrightarrow \pi(a,c)
\]
on arrows. It is immediate that it is associative (the composition of paths being so), and unital. It admits as identity arrow at $x$ either the ``empty loop at $x$'', which is (directly) homotopic in $\Ch(x,x)$ to the ``degenerate'' path $x \xrightarrow{\mathrm{id}} x$
\msp{page 62}
of type $\Delta_+$, or also to the degenerate path of any type $\tau$, say
\[
x \xleftarrow{\mathrm{id}} x \xrightarrow{\mathrm{id}} x \xrightarrow{\mathrm{id}} x \xleftarrow{\mathrm{id}} x \xleftarrow{\mathrm{id}} x,
\]
such a path being none other than $p_\tau^*(x)$, where
\[
p_\tau : \tau \longrightarrow 1
\]
is the degeneracy morphism (the unique morphism in $\Ty$ to 1, which is a final object), whence
\[
p_\tau^* : \Ch_1(x,x) = \{x\} \longrightarrow \Ch_\tau(x,x) \subseteq \Ch(x,x).
\]
The category $\Pi(X)$ thus obtained must be called the \textit{fundamental groupoid} of $X$. Indeed, % [sic] obtenu -> obtenue, agreement with « La catégorie »

\begin{enumerate}[label=\alph*)]
\item all the arrows of $\Pi(X)$ are invertible, i.e. $\Pi(X)$ is indeed a groupoid, and
\item the canonical functor
\[
X \longrightarrow \Pi(X)
\]
which is the identity on objects, and which is defined on arrows by $u \longmapsto c(u)$ as above, is \textit{universal} (in Cat) for functors (i.e. Cat-morphisms) from $X$ to \textit{groupoids.} (More generally still, it is universal for functors $X \xrightarrow{F} Y$ such that $F$ transforms every arrow of $X$ into \texttt{[an]} invertible arrow.)
\end{enumerate}
% foot: Version 22 décembre 2001

\printedpage{42}
% head: A. Grothendieck : Dérivateurs, Ch. VII, Édition des Universités de Montpellier II et Paris VII
\msp{page 63}
To prove that every arrow in $\Pi(X)$ is invertible, since every nonempty path is a composite of ``elementary'' paths of type $\Delta_+$ or $\Delta_-$, it suffices to prove that these paths are invertible. In fact, they correspond bijectively,
\[
\Ch_{\Delta_+}(a,b) \simeq \Ch_{\Delta_-}(b,a),
\]
by associating to the path from $a$ to $b$
\[
c(u) : a \xrightarrow{u} b \qquad \text{of type } \Delta_+
\]
the path
\[
c(u)^o : b \xleftarrow{u} a \qquad \text{of type } \Delta_-
\]
from $b$ to $a$. I say that
\[
\begin{array}{rcl}
c(u)^o c(u) & \sim & 1_a\\
c(u) c(u)^o & \sim & 1_b
\end{array}
\]
\[
\left.
\begin{tikzcd}
 & a \arrow[r, "u"] & b & a \arrow[l, "u"'] \\
 {} & a \arrow[u, "\mathrm{id}"'] \arrow[r, "\mathrm{id}"] & a \arrow[u, "u"'] & a \arrow[u, "\mathrm{id}"'] \arrow[l, "\mathrm{id}"'] \\
 {} \arrow[u] & & a &
\end{tikzcd}
\right\}\
\begin{array}{l}
\text{we have a homomorphism in } \Ch_{\Delta_+\Delta_-}(a,a)\\
\text{from } 1_a \text{ to } c(u)^o c(u).
\end{array}
\]
\[
\left.
\begin{tikzcd}
 & b & a \arrow[l, "u"'] \arrow[r, "u"] & b \\
 {} & a \arrow[u, "u"'] & a \arrow[u, "\mathrm{id}"'] \arrow[l, "\mathrm{id}"'] \arrow[r, "\mathrm{id}"] & a \arrow[u, "u"'] \\
 {} \arrow[u] & & a &
\end{tikzcd}
\right\}\
\begin{array}{l}
\text{likewise in } \Ch_{\Delta_-\Delta_+}(b,b),\\
\text{from } 1_b \longrightarrow c(u)c(u)^0 % [sic] c(u)^0 -> c(u)^o, digit 0 printed instead of the letter o used elsewhere
\end{array}
\]
\texttt{[rather} $1_b \longleftarrow c(u)c(u)^0$\texttt{, viz.} % [sic] c(u)^0 -> c(u)^o, digit 0 printed instead of the letter o used elsewhere
\[
\begin{tikzcd}
b \arrow[d, "\mathrm{id}"] & a \arrow[l, "u"'] \arrow[r, "u"] \arrow[d, "u"] & b \arrow[d, "\mathrm{id}"] \\
b & b \arrow[l, "\mathrm{id}"'] \arrow[r, "\mathrm{id}"] & b \ \texttt{].}
\end{tikzcd}
\]
% foot: Version 22 décembre 2001

\printedpage{43}
% head: A. Grothendieck : Dérivateurs, Ch. VII, Édition des Universités de Montpellier II et Paris VII
\msp{page 64}
Let us note that, more generally, for any path
\[
c : a \longrightarrow b, \qquad \text{of type } \tau,
\]
we define an \textit{opposite path}
\[
c^o : b \longrightarrow a \qquad \text{of type } \tau^o,
\]
where
\[
\tau \longmapsto \tau^o : \mathrm{Ty} \longrightarrow \mathrm{Ty}
\]
is the anti-involution of the free monoid Ty which exchanges the two generators $\Delta_-$, $\Delta_+$, and which is therefore given on the general element by
\[
(\Delta_{\varepsilon_1}\Delta_{\varepsilon_2}\cdots\Delta_{\varepsilon_n})^o = \Delta_{-\varepsilon_n}\Delta_{-\varepsilon_{n-1}}\cdots\Delta_{-\varepsilon_1}.
\]
By induction on the length of $c$ (\footnote[18]{unnecessary, follows directly from formula $(*)$ \texttt{[on the]} following page, by decomposing $c$ into composites of ``elementary'' paths of type $\Delta_+$ or $\Delta_-$.}), using the preceding result, that $c(u)$, $c(u)^o$ are inverses of one another in $\Pi(X)$, we find
\[
\left.
\begin{array}{rcl}
c^o c & = & \mathrm{id}_a\\
c c^o & = & \mathrm{id}_b
\end{array}
\right\}\text{in } \Pi(X).
\]

\textbf{Example.} If
\[
\begin{array}{rcl}
\tau & = & \Delta_+\Delta_-\Delta_+^2\Delta_-\\
\tau^o & = & \Delta_+\Delta_-^2\Delta_+\Delta_-,
\end{array}
\]
and if we have a path
\[
c : a_0 \xrightarrow{u_1} a_1 \xleftarrow{u_2} a_2 \xrightarrow{u_3} a_3 \xrightarrow{u_4} a_4 \xleftarrow{u_5} a_5
\]
of type $\tau$, we deduce from it the path
\[
c^o : a_5 \xrightarrow{u_5} a_4 \xleftarrow{u_4} a_3 \xleftarrow{u_3} a_2 \xrightarrow{u_2} a_1 \xleftarrow{u_1} a_0
\]
of type $\tau^o$. The operation $c \longmapsto c^o$ satisfies in
\msp{page 65}
fact \texttt{[for]}
\[
\begin{array}{ccccc}
a & \xrightarrow{\gamma} & b & \xrightarrow{\gamma'} & c\\
c & \xrightarrow{\gamma'^o} & b & \xrightarrow{\gamma^o} & a
\end{array}
\]
the relation
\[
(\gamma'\gamma)^o = \gamma^o\gamma'^o, \tag*{$(*)$}
\]
% foot: Version 22 décembre 2001